% ===========================================================================
% Section 1 -- Introduction: Can Machines Invent?
% File: sections/01_introduction.tex
% Reader question: Can an AI create knowledge no human wrote down?
% One-sentence answer: Skeptics say models trained on history can only
%   recombine it. AlphaGo Zero showed that a learner with a referee can go
%   beyond history, and we test that for invention.
% Core points advanced: 4 (skeptics right about imitation, wrong about
%   learning), 3 (search = 1), 1 (innovation is a game).
% Exhibits: Figure 1 (front matter, label fig:emptysquare) and Box 1.
% Budget: 1.5 pages. Editor notes: sections/01_introduction.notes.md
% ===========================================================================

\section{Can Machines Invent?}\label{sec:intro}

Can a machine create knowledge that no human has written down? The strongest
skeptics say no: a model trained on the record of human knowledge can only
recombine it. Go suggests another answer. AlphaGo Zero, taught by a referee
instead of a record, went beyond the best human play without seeing a human
game \citep{silver2017mastering}. We ask whether a similar training signal
can improve algorithm invention: can learning from past evaluations make
the first proposal on a new problem more useful?

\paragraph{The skeptics' challenge.}
In March 2025, Yann LeCun argued that language models retrieve existing
solutions rather than find new ones: ``in their purest form, they are incapable
of inventing new things'' \citep{lecun2025bigtech}. The same month, Thomas Wolf
warned that current trends will give us ``a country of yes-men on servers''
\citep{wolf2025einstein}. Dwarkesh Patel had asked Dario Amodei in August 2023
what to make of models with ``basically the entire corpus of human knowledge
memorized'' that had not made ``a single new connection that has led to a
discovery'' \citep{patel2023question}. In January 2025, Demis Hassabis posed
the hardest version as a benchmark for general intelligence: today's systems
``can play a game of Go at a world champion level. But could a system invent
Go?'' \citep{hassabis2025bigtech}. These challenges separate producing an
unfamiliar idea from learning which ideas will work. Predicting historical
text does not directly reward improvement over an incumbent. That motivates
a different training signal; it does not establish a ceiling on what a
pretrained model can discover (\cref{sec:anatomy}).

\paragraph{Humans are bounded by history too.}
Research effort keeps rising while its productivity
falls: Moore's law now needs more than 18 times as many researchers as in the early 1970s \citep{bloom2020ideas}. Under a growing burden
of knowledge, inventors reach the frontier later and specialize more narrowly
\citep{jones2009burden}. Papers and patents have grown less likely to
break with the past \citep{park2023disruptive}, and fields open to outsiders
when their leading scientists die \citep{azoulay2019funeral}. The bottleneck is
not a shortage of minds. It is the gravity of what we already know.

\paragraph{From search to learning.}
Automated research systems already couple proposals to tests. The AI
Scientist \citep{lu2024aiscientist} and the AI co-scientist
\citep{gottweis2025coscientist} run agentic research loops; AlphaEvolve evolves
programs against automated evaluators \citep{novikov2025alphaevolve}; and
Karpathy's autoresearch retains training changes that improve a measured
objective \citep{karpathy2026autoresearch}. Socher's \emph{Eureka Machine}
extends this vision to agents working with simulators and autonomous
laboratories \citep{socher2026bookbite,socher2026eureka}. These systems make
judgment executable. Our question is what their evaluations can teach the
system that proposes the next idea. Keeping a winning program preserves an
output; training a proposer on those outcomes may preserve a capability that
transfers to a new problem.

\begin{keyidea}
Can judgment bought during training improve the first proposal at deployment?
\end{keyidea}

\paragraph{The Zero move.}
Go faced the same choice and took the other road. AlphaGo learned from human
expert games and searched heavily at every move \citep{silver2016mastering}.
AlphaGo Zero learned only from games against itself, and its network alone,
with no search, rated 3,055 Elo. That is close to the 3,144 of AlphaGo Fan, the
version that beat the European champion Fan Hui
\citep{silver2016mastering,silver2017mastering}. This motivates an analogous
goal for invention: learn from repeated program evaluations during training,
then deploy a writer with a single candidate budget, \searchone{}
(\hyperref[box:searchone]{Box~1}). The budget is a test condition, not a
promise that every first proposal wins. \Cref{fig:emptysquare} summarizes
the distinction between spending search at deployment and learning from it
in advance.

\begin{explainer}{Box 1: What ``search = 1'' means}
\phantomsection\label{box:searchone}%
\begin{itemize}[leftmargin=*,nosep]
  \item \textbf{One inference, one candidate.} For a new problem, the program
    writer emits one complete program ($B = 1$), with no retries, repairs,
    best-of-$N$, reranking or weight updates.
  \item \textbf{Everything is counted,} including reasoning tokens and problem
    encoding. One inference is one generation, not one forward pass.
  \item \textbf{Declare the fitting procedure.} Fitting the emitted algorithm
    to training rows is allowed. Any internal model or hyperparameter
    selection must be disclosed and counted in its compute budget. One
    outer proposal does not establish the absence of internal search.
    Training search is priced separately in \cref{sec:price}.
\end{itemize}
\textbf{A single candidate, with prior knowledge:} the writer starts from
pretraining and human ideas. The Zero analogy concerns the training signal;
it does not imply that the system starts without human knowledge.
\end{explainer}

\begin{figure}[t]
  \centering
  \fbox{\parbox{0.88\linewidth}{\TBD{Conceptual diagram: distinguish search
  during training from the candidate budget at deployment.}}}
  \caption{\textbf{Learning before deployment and searching during it.}
  A conceptual comparison is pending; no measured \EZ{} advantage is implied.}
  \label{fig:emptysquare}
\end{figure}

\paragraph{The proposal and its test.}
We formulate algorithm invention as a game: the encoded problem is the board,
a complete program is the move, a verifier is the referee and the incumbent
algorithm is the opponent (\cref{sec:game}). \EZ{} combines a program writer
with a problem--program map intended to guide learning from those contests
(\cref{sec:method}). The central hypothesis is that this training improves
first proposals on unseen problems. The current learning record is reported
separately from the proposed deployment comparisons.

\Cref{sec:results} specifies the evidence needed to test that hypothesis:
paired comparisons against algorithms proposed by Claude, Codex and Gemini;
a comparison with KumoRFM \citep{kumorfm2025}; a performance curve against
agents allowed repeated attempts; and an ablation of the problem--program
map. \TBD{Complete these comparisons and document task coverage, model
versions, prompts, budgets and uncertainty before claiming an advantage.}
The paper develops five elements:
\begin{enumerate}[leftmargin=*,nosep,label=\arabic*.]
  \item \textbf{The anatomy of innovation:} four pillars (aim, variation,
    judgment and retention), with an idealized relation between proposal
    quality and the number of attempts (\cref{sec:anatomy}).
  \item \textbf{The innovation game:} rules that make invention playable on
    many boards (\cref{sec:game}).
  \item \textbf{The economics of algorithmic self-play:} an accounting of
    training cost and the cost per verified improvement (\cref{sec:price}).
  \item \textbf{\EZ{}:} a writer and a problem--program map designed to learn
    from referee feedback (\cref{sec:method}).
  \item \textbf{An evaluation plan:} tests of first-proposal quality,
    deployment search and the map's contribution (\cref{sec:results}).
\end{enumerate}
\Cref{sec:road} traces each pillar's history in AI; \cref{sec:meaning} asks
what changes and where the claim stops.

To see what the referee can teach, we first separate the parts of innovation.
